\section{Preliminaries}

\subsection{Problem Formulation}

We study mixed-autonomy dynamic routing: connected and autonomous vehicles
(CAVs) and human-driven vehicles (HDVs) travel over a directed road network
under time-varying departure demand, and a controller routes the CAVs.

\noindent\textbf{Definition 1 (Segment-Node Road Network).}
The road network is a directed graph \(G=(N,E)\): each node \(u\in N\) is a
road segment with static attributes (length, speed limit, lane count,
capacity, free-flow travel time), and a directed edge \((u,v)\in E\) exists
when vehicles can move directly from segment \(u\) to \(v\).  We adopt this
segment-node representation---nodes are road segments rather than
intersections, and edges represent allowed movements between consecutive
segments---so that traffic quantities and route actions share one
coordinate system.  A fixed node ordering assigns each segment a stable
index, with adjacency matrix \(A\).

\noindent\textbf{Definition 2 (Vehicles).}
The departure horizon is partitioned into decision epochs
\(t=0,\ldots,T\), with active population
\(\mathcal{N}_t=\mathcal{N}^{\mathrm{CAV}}_t\cup\mathcal{N}^{\mathrm{HDV}}_t\)
at epoch \(t\).  Each vehicle \(i\) has an origin segment \(o_i\), a
destination segment \(g_i\), a scheduled departure time
\(\tau_i^{\mathrm{dep}}\), and a type in \(\{\mathrm{CAV},\mathrm{HDV}\}\);
\(u_{i,t}\) denotes its current segment when active.  CAVs are controlled by
the routing policy; HDVs are self-interested and cannot be rerouted.

\noindent\textbf{Definition 3 (Route Candidates and Decision Timing).}
An active CAV \(i\) at epoch \(t\) selects one route from its candidates
\begin{equation}
\mathcal{P}_{i,t}
=
\{P^1_{i,t},\ldots,P^{K}_{i,t}\},
\end{equation}
where each candidate is a segment sequence from \(u_{i,t}\) to \(g_i\),
generated by repeated Dijkstra search with a route-diversity penalty
(Appendix~\ref{app:math-details}), and the action is the route index
\(a_{i,t}\in\{1,\ldots,K\}\).  An HDV chooses its route only at its
departure time \(\tau_i^{\mathrm{dep}}\).

\subsection{Markov Decision Process}

We model the route-choice process as a Markov decision process
\(\mathcal{M}=(\mathcal{S},\mathcal{A},\mathcal{T},r,\gamma)\) over the
global state.  Each CAV, however, acts from its own observation
\(o_{i,t}\) rather than the full \(s_t\), so execution is partially
observable; the learned components introduced below are trained with access
to the global state, following centralized training with decentralized
execution.

\noindent\textbf{Definition 4 (Observation/State).}
The global state \(s_t\in\mathcal{S}\) includes the road graph, the
scheduled vehicles with their attributes, and the current route candidates.
Each active CAV observes a projection \(o_{i,t}=\mathcal{O}_i(s_t)\)
containing its current position, destination, candidate routes, and
route-level traffic summaries.

\noindent\textbf{Definition 5 (Action).}
Let \(\mathcal{I}_t=\mathcal{N}^{\mathrm{CAV}}_t\) denote the active CAVs at
epoch \(t\).  The joint action is
\begin{equation}
\begin{array}{rcl}
a_t&=&(a_{i,t})_{i\in\mathcal{I}_t},
\end{array}
\end{equation}
where \(a_{i,t}=k\) selects candidate route \(P^k_{i,t}\).  HDV route
choices are generated by background human-routing behavior and are part of
the environment.

\noindent\textbf{Definition 6 (Transition).}
Given the current state, the joint CAV action, the departure demand \(D_t\)
of the vehicles scheduled to depart, and the HDV traffic state
\(H_t^{\mathrm{HDV}}\) induced by the HDVs' own route choices, the
transition is
\begin{equation}
\begin{array}{l}
s_{t+1}\sim
\mathcal{T}(\cdot\mid s_t,a_t,D_t,H_t^{\mathrm{HDV}}).
\end{array}
\end{equation}
It injects newly departing vehicles, applies the selected routes, advances
all vehicles, and updates segment speeds, occupancies, and active sets.

\noindent\textbf{Definition 7 (Objective/Reward).}
Each completed vehicle \(i\) (CAV or HDV) has travel time
\(\Delta_i=t_i^{\mathrm{arr}}-\tau_i^{\mathrm{dep}}\).  The objective is
system-level: minimize the expected global mean travel time over all
completed vehicles,
\begin{equation}
\begin{array}{rcl}
\min\quad
\mathcal{T}_{\mathrm{all}}&=&
\displaystyle
{E}_{\pi_\theta}
\left[
\frac{1}{|\mathcal{N}^{\mathrm{arr}}|}
\sum_{i\in\mathcal{N}^{\mathrm{arr}}}
\Delta_i
\right],
\end{array}
\end{equation}
where \(\mathcal{N}^{\mathrm{arr}}\) is the set of arrived vehicles, and the
per-step rewards \(r\) used by the learner are instantiated in
Section~\ref{sec:method}.

\subsection{World-Model Formulation}
\label{sec:wm-formulation}

The transition \(\mathcal{T}\) is accessible only through microscopic
simulation: evaluating a joint action requires rolling out an urban-scale
episode, and the network-wide effect of a route choice emerges only after
many decision intervals, so direct policy optimization on \(\mathcal{M}\)
consumes a large number of costly simulator episodes.  We instead learn a
model of the routing process and optimize the policy primarily inside it.
The transition \(\mathcal{T}\) couples two processes: the evolution of the
network-wide traffic field and the progress of individual vehicles along
their routes.  We model the two at separate levels---a learned macroscopic
traffic transition and an analytic microscopic vehicle advancement---and
compose them into the imagined rollouts used for policy learning.

\paragraph{Macroscopic traffic level.}
The network-wide state is the segment traffic field \(X_t=\chi(s_t,G)\):
\begin{equation}
\begin{array}{l}
X_t(u)=
[\mathrm{c}_t(u),\ v_t(u),
\kappa(u),
\ell(u),\ v_{\max}(u)],
\end{array}
\label{eq:traffic-field}
\end{equation}
with vehicle count \(\mathrm{c}_t\), speed \(v_t\), capacity \(\kappa\),
length \(\ell\), and speed limit \(v_{\max}\).  \(X_t\) is a lossy
segment-aligned view of \(s_t\) that removes vehicle identities,
destinations, and chosen route suffixes; these components re-enter through
the microscopic level below.  The macroscopic action is the planned CAV
pressure field \(\Psi_t^{\mathrm{CAV}}\), a segment-level distribution of
where the controlled vehicles intend to travel, obtained by projecting the
joint route choices onto the same segment graph.  The traffic field is
encoded as \(Z_t=f_\phi(X_t,G)\), and a learned transition \(F_\phi\),
both parameterized by neural networks detailed in
Section~\ref{sec:method}, predicts the next latent
traffic state under the macroscopic action and the exogenous conditions:
\begin{equation}
  \begin{array}{rcl}
  \hat Z_{t+1}
  &=&
  F_\phi(
  Z_t; \Psi_t^{\mathrm{CAV}},D_t,H_t^{\mathrm{HDV}}
  ),
  \end{array}
\end{equation}

\paragraph{Microscopic vehicle level.}
Each active CAV is tracked individually: its state is the observation
\(o_{i,t}\), together with its current segment, route suffix, and arrival
status, and its action is the route index \(a_{i,t}\).  A deterministic
vehicle wrapper advances each vehicle along its chosen route by consuming
the decoded segment travel times within one decision interval:
\begin{equation}
(u_{i,t+1}, P^{\mathrm{rem}}_{i,t+1})
=
\mathrm{Advance}\!\big(u_{i,t}, P^{a_{i,t}}_{i,t};
\widehat\tau_{t+1}, \Delta t\big),
\qquad
\widehat\tau_t(u) = \frac{\ell(u)}{\widehat v_{t,u}},
\label{eq:micro-advance}
\end{equation}
where \(\widehat v_{t,u}\) is the segment speed decoded from the
macroscopic state, vehicles reaching their departure time are added, and
arrived vehicles are removed to form the next active set.  Vehicle progress
and the traffic field thus evolve from the same prediction.

\paragraph{Rewards at the two levels.}
The reward signals mirror the two levels.  The microscopic per-vehicle
reward charges each CAV the imagined travel time it accumulates before its
next decision or arrival, providing vehicle-specific credit.  The
macroscopic global reward charges the travel time accumulated by all
vehicles present in each decision interval, so that its episode-level sum
matches the objective of Definition~7.  Both rewards are computed from
decoded readouts rather than learned; their exact forms are given in
Section~\ref{sec:method}.

The two levels are coupled in both directions: microscopic route choices
are aggregated into the macroscopic pressure that conditions the
transition, and macroscopic predictions supply the travel times that drive
microscopic progress.  Imagined rollouts at the two levels supply the
transitions for actor-critic updates, while simulator episodes remain the
source of model fitting and final policy evaluation.
